\documentclass[11pt,a4paper]{report}
\usepackage[T1]{fontenc}    % Indispensable pour la coupure des mots en français
\usepackage[utf8]{inputenc} 
\usepackage{amsmath, amsthm}
\usepackage{amssymb}
\usepackage{amsfonts}
\usepackage{stmaryrd}
\usepackage{xcolor}
\usepackage{titlesec}
\titleformat{\chapter}[block]
  {\normalfont\bfseries\filcenter}
  {}
  {0pt}
  {\titlerule\vspace{1ex}\Huge}
\usepackage[margin=1cm, includeheadfoot]{geometry}
\usepackage{frenchmath}
% --- LE SECRET EST ICI : Le package tcolorbox ---
\usepackage[most]{tcolorbox}

% On crée un environnement "commentaire" personnalisé avec la barre grise
\newtcolorbox{commentaire}{
  blanker,                  % Supprime le fond et les bordures par défaut
  left=15pt,                % Espace entre la barre et le texte
  top=3pt,
  bottom=3pt,
  before skip=10pt,         % Espace avant le bloc
  after skip=10pt,          % Espace après le bloc
  borderline west={3pt}{0pt}{gray!40}, % Crée la barre verticale (épaisseur 3pt, couleur gris clair)
  fontupper=\itshape\color{black!70}   % Tout le texte à l'intérieur en italique et gris foncé
}
\begin{document}

\chapter{Mines-Ponts}

% --- TITRE DE L'EXERCICE (Style identique au titre 1.7 de l'image) ---
\begin{center}
    \large \textbf{1. Etude d'une matrice à diagonale nulle}
\end{center}
\vspace{0.2cm}

% --- ENCADRÉ DE L'ÉNONCÉ (Même style visuel avec source en bas à droite) ---
\noindent\framebox[\textwidth]{%
    \begin{minipage}{0.94\textwidth}
        \vspace{0.2cm}
        Soit un entier $n \geqslant 2$. On note $A$ la matrice de $\mathcal{M}_n(\mathbb{R})$ dont les coefficients diagonaux valent 0 et les autres coefficients valent 1.
        \vspace{0.2cm}
        
        \noindent $1.$ Calculer $A^2$. En déduire que $A$ est inversible et exprimer $A^{-1}$ ;
        
        \noindent $2.$ Déterminer les valeurs propres de $A$ et les espaces propres correspondants.
        \vspace{0.1cm}
        \begin{flushright}
            (\textbf{Mines-Ponts, PC})
        \end{flushright}
    \end{minipage}%
}
\vspace{0.4cm}

% --- DÉBUT DE LA SOLUTION (Introduction avec triangle \triangleright) ---
\noindent $\triangleright$ \textbf{Solution.} \\
% --- PREMIÈRE QUESTION (Introduite par une puce \bullet comme dans l'image) ---
\textbf{1.} Commençons par calculer $A^2$. Soit $(i,j) \in \llbracket 1,n \rrbracket^2$. On a :
\[
[A^2]_{i,j} = \sum_{k=1}^n A_{i,k} A_{k,j} = \sum_{\substack{k=1 \\ k \notin \{i,j\}}}^n A_{i,k} A_{k,j} = \sum_{\substack{k=1 \\ k \notin \{i,j\}}}^n 1 = \left| \llbracket 1,n \rrbracket \setminus \{i,j\} \right|,
\]
ce qui donne:
\[
[A^2]_{i,j} = \begin{cases} n-2 & \text{si } i \neq j \\ n-1 & \text{si } i = j \end{cases}
\]
On en déduit que 
\[
A^2 = (n-2)A + (n-1)I_n.
\]
On peut réécrire cette relation sous la forme $A \cdot \left( \frac{1}{n-1}(A - (n-2)I_n) \right) = I_n$. Ainsi, \fbox{la matrice $A$ est inversible} et, par unicité de l'inverse, on obtient :
\[
\boxed{A^{-1} = \frac{1}{n-1}(A - (n-2)I_n).}
\]

\vspace{0.4cm}

% --- DEUXIÈME QUESTION (Deuxième puce \bullet) ---
\textbf{2.} Déterminons les valeurs propres de $A$. L'équation aux valeurs propres $AX = \lambda X$, avec $\lambda \in \mathbb{C}$ et $X = (x_1, \dots, x_n)^T \in \mathbb{R}^n \setminus \{0\}$, équivaut au système suivant :
\[
\forall i \in \llbracket 1,n \rrbracket, \quad \sum_{\substack{j=1 \\ j \neq i}}^n x_j = \lambda x_i \quad (1)
\]
En sommant ces $n$ identités, il vient :
\[
(n-1) \left( \sum_{j=1}^n x_j \right) = \lambda \left( \sum_{i=1}^n x_i \right).
\]
Deux cas se présentent alors:
\begin{itemize}
    \item[\textbf{--}] Si $\sum_{i=1}^n x_i = 0$ : le vecteur $X$ appartient à l'hyperplan $H = \left\{ X = (x_1, \dots, x_n)^T \in \mathbb{R}^n \;\left|\; \sum_{i=1}^n x_i = 0 \right.\right\}$. On déduit alors de la relation (1) que :
    \[
    \sum_{j=1}^n x_j - x_i = \lambda x_i \quad \forall i \in \llbracket 1,n \rrbracket,
    \]
    soit $-x_i = \lambda x_i$, d'où $\lambda = -1$.
    \item[\textbf{--}] Si $\sum_{i=1}^n x_i \neq 0$ : on obtient directement $\lambda = n-1$.
\end{itemize}
On peut donc déduire que le spectre de $A$ est :
\[
\boxed{\mathrm{Sp}(A) = \{n-1, -1\},}
\]
et les espaces propres associés sont donnés par :
\[
\boxed{E_{n-1}(A) = \mathrm{Vect}\left((1, \dots, 1)^T\right) \quad \text{et} \quad E_{-1}(A) = H.}
\]

% --- TITRE DE L'EXERCICE ---
\begin{center}
    \large \textbf{2. Suites de Cauchy et complétude}
\end{center}
\vspace{0.2cm}

% --- ENCADRÉ DE L'ÉNONCÉ ---
\noindent\framebox[\textwidth]{%
    \begin{minipage}{0.94\textwidth}
        \vspace{0.2cm}
        Soit $(E, \| \cdot \|)$ un espace vectoriel normé. Une suite $(u_n) \in E^{\mathbb{N}}$ est de Cauchy si $\forall \varepsilon > 0, \exists N \in \mathbb{N}, \forall n, m \geqslant N, \|u_n - u_m\| \leqslant \varepsilon$.
        \vspace{0.2cm}
        
        \noindent \textbf{a)} Montrer que toute suite convergente est de Cauchy.
        
        \noindent \textbf{b)} Dans $E = \mathbb{R}[X]$ muni de la norme $\left\| \sum a_k X^k \right\| = \max |a_k|$, montrer que la suite $(P_n)$ de terme général $P_n = 1 + \sum_{k=1}^n \frac{X^k}{k}$ est de Cauchy sans être convergente ;
        
        \noindent \textbf{c)} Montrer que toute suite de Cauchy est bornée.
        
        \noindent \textbf{d)} Montrer que, si $(u_n)$ est de Cauchy et possède une suite extraite convergente, alors $(u_n)$ est convergente.
        
        \noindent \textbf{e)} On admet le théorème de Bolzano-Weierstrass dans $\mathbb{R}$. Montrer que si $E$ est de dimension finie, alors la suite $(u_n)$ est convergente si et seulement si elle est de Cauchy.
        \vspace{0.1cm}
        \begin{flushright}
            (\textbf{Mines-Ponts, PSI (2025)})
        \end{flushright}
    \end{minipage}%
}
\vspace{0.4cm}

\noindent $\triangleright$ \textbf{Solution.} \\
\noindent \textbf{a)} Soit $(u_n)_{n \in \mathbb{N}}$ une suite convergente. Notons $l = \lim_{n \to +\infty} u_n$. Soit $\varepsilon > 0$. Il existe $N \in \mathbb{N}$ tel que pour tout $n \geqslant N$, $\|u_n - l\| \leqslant \varepsilon/2$.

\noindent Pour $m \geqslant N$, on a également $\|u_m - l\| \leqslant \varepsilon/2$. Ainsi, l'inégalité triangulaire donne :
$$ \|u_n - u_m\| \leqslant \|u_n - l\| + \|u_m - l\| \leqslant \varepsilon/2 + \varepsilon/2 = \varepsilon $$
\noindent d'où :
\begin{center}
    \fbox{$(u_n)_{n \in \mathbb{N}}$ est de Cauchy}
\end{center}

\vspace{0.5cm}

\noindent \textbf{b)} Montrons d'abord que pour la norme définie, $(P_n)_{n \in \mathbb{N}}$ n'est pas convergente. On raisonne par l'absurde en supposant qu'elle est convergente. On note $P = \lim_{\|\cdot\| \ n \to +\infty} P_n = \sum_{k=0}^d a_k X^k$, avec $(a_k)_{0 \leqslant k \leqslant d} \in \mathbb{R}^{d+1}$ et $d \in \mathbb{N}$.

\noindent On a donc :
$$ \|P_n - P\| = \left\| 1 + \sum_{k=1}^n \frac{X^k}{k} - \sum_{k=0}^d a_k X^k \right\| $$
\noindent D'où, pour $n > d$ : 
$$ \|P_n - P\| = \max \left( |1-a_0|, \left| \frac{1}{k} - a_k \right| _{1 \leqslant k \leqslant d}, \left( \frac{1}{k} \right)_{d+1 \leqslant k \leqslant n} \right) \geqslant \frac{1}{d+1} \neq 0 $$
\noindent On déduit que :
\begin{center}
    \fbox{\begin{tabular}{c} $(P_n)_{n \in \mathbb{N}}$ n'est pas convergente \\ pour la norme $\|\cdot\|$ \end{tabular}}
\end{center}

\noindent $\bullet$ Montrons que $(P_n)_{n \in \mathbb{N}}$ est néanmoins de Cauchy. Soit $\varepsilon > 0$. Remarquons d'abord que pour $m > n$ :
$$ P_m - P_n = \sum_{k=n+1}^m \frac{X^k}{k} $$
\noindent d'où :
$$ \|P_m - P_n\| = \max_{n+1 \leqslant k \leqslant m} \left( \frac{1}{k} \right) $$

\noindent Puisque $\frac{1}{k} \leqslant \frac{1}{n+1}$ pour tout $k \in \llbracket n+1, m \rrbracket$, notons $N = \left\lceil \frac{1}{\varepsilon} - 1 \right\rceil$. On a donc pour $m > n \geqslant N$ :
$$ \|P_m - P_n\| \leqslant \varepsilon $$
\noindent d'où :
\begin{center}
    \fbox{$(P_n)_{n \in \mathbb{N}}$ est de Cauchy}
\end{center}

\vspace{0.5cm}

\noindent \textbf{c)} Soit $(u_n)_{n \in \mathbb{N}}$ une suite de Cauchy. Il existe $N \in \mathbb{N}$ tel que pour tous $n,m \geqslant N$, $\|u_n - u_m\| \leqslant 1$. Pour $m = N$, l'inégalité triangulaire renversée donne : $\|u_n\| - \|u_N\| \leqslant \|u_n - u_N\|$, donc pour tout $n \geqslant N$ :
$$ \boxed{ \|u_n\| \leqslant \|u_N\| + 1 } $$

\noindent $(u_n)_{n \in \mathbb{N}}$ est donc bornée à partir d'un certain rang et on déduit que :
\begin{center}
    \fbox{$(u_n)_{n \in \mathbb{N}}$ est bornée}
\end{center}

\vspace{0.5cm}

\noindent \textbf{d)} Soit $(u_n)_{n \in \mathbb{N}}$ une suite de Cauchy qui possède une suite extraite convergente $(u_{\varphi(n)})_{n \in \mathbb{N}}$ de limite $l \in E$.

\noindent Soit $\varepsilon > 0$. Il existe, puisque $(u_{\varphi(n)})_{n \in \mathbb{N}}$ est convergente, un entier $N_0 \in \mathbb{N}$ tel que pour tout $n \geqslant N_0$, $\|u_{\varphi(n)} - l\| \leqslant \varepsilon/2$.

\noindent Puisque $(u_n)_{n \in \mathbb{N}}$ est de Cauchy, il existe $N_1 \in \mathbb{N}$ tel que pour tous $n,m \geqslant N_1$, $\|u_n - u_m\| \leqslant \varepsilon/2$. On choisit $m = \varphi(K)$ avec $K = \max(N_0,N_1)$. On a donc pour tout $n \geqslant K$ :
$$ \|u_n - l\| \leqslant \|u_n - u_m\| + \|u_m - l\| \leqslant \varepsilon/2 + \varepsilon/2 = \varepsilon $$
\noindent d'où :
\begin{center}
    \fbox{$(u_n)_{n \in \mathbb{N}}$ est convergente.}
\end{center}

\vspace{0.5cm}

\noindent \textbf{e)} Rappelons le théorème de Bolzano-Weierstrass, hors-programme pour les élèves PCSI/PSI* :
\begin{center}
    \fbox{Toute suite de réels bornée admet une sous-suite convergente.}
\end{center}

\noindent $\bullet$ Le sens "$\Rightarrow$" a été déjà démontré en \textbf{a)}.

\noindent $\bullet$ Montrons le sens "$\Leftarrow$" : Soit $(u_m)_{m \in \mathbb{N}} = \left( \sum_{k=1}^n u_m^{(k)} e_k \right)_{m \in \mathbb{N}}$ une suite dans $E$ (où $n = \dim(E)$ et $(e_1, \ldots, e_n)$ est une base de $E$).

\noindent Supposons que $(u_m)_{m \in \mathbb{N}}$ est de Cauchy. Elle est donc bornée (d'après \textbf{c)}). Qu'importe la norme $\|\cdot\|$, l'équivalence des normes en dimension finie permet d'affirmer que les suites coordonnées $(u_m^{(k)})_{m \in \mathbb{N}}$ pour $k \in \llbracket 1, n \rrbracket$ sont bornées. (On compare avec la norme $\|u\|_\infty = \max_{1 \leqslant k \leqslant n} (\|u^{(k)}\|_\infty)$, où $\|u^{(k)}\|_\infty = \sup_{m \in \mathbb{N}} |u_m^{(k)}|$ est une norme sur l'espace des suites bornées).

\noindent D'après le théorème admis, les suites $(u_m^{(k)})_{m \in \mathbb{N}}$ admettent des sous-suites convergentes. En imbriquant les extractrices de chaque sous-suite, on peut ainsi vérifier que $(u_m)_{m \in \mathbb{N}}$ possède une sous-suite convergente. D'où, par \textbf{d)} :
\begin{center}
    \fbox{$(u_m)_{m \in \mathbb{N}}$ est convergente.}
\end{center} 

\pagebreak

% --- TITRE DE L'EXERCICE ---
\begin{center}
    \large \textbf{3. Etude d'une relation d'anticommutation}
\end{center}

\vspace{0.2cm}
% --- ENCADRÉ DE L'ÉNONCÉ ---
\noindent\framebox[\textwidth]{%
\begin{minipage}{0.94\textwidth}
    \vspace{0.2cm}
    Soit $E$ un $\mathbb{K}$-espace vectoriel de dimension $n$. Soient $U$ et $V$ deux sous-espaces vectoriels de $\mathcal{L}(E)$ tels que $U + V = \mathcal{L}(E)$ et tels que $\forall u \in U, \forall v \in V, u \circ v + v \circ u = 0$.
    
    \vspace{0.2cm}
    \noindent \textbf{a)} Montrer qu'il existe $p \in U$ et $q \in V$ deux projecteurs tels que $p + q = \mathrm{id}$.
    
    \noindent \textbf{b)} Si $v \in V$, on pose $\varphi(v) = v_{|\operatorname{Ker}(p)} \in \mathcal{L}(\operatorname{Ker} p, E)$. Montrer que $\varphi$ est injective. En déduire que $\dim(V) \leqslant (n - r)^2$, avec $r = \operatorname{rg}(p)$.
    
    \noindent \textbf{c)} Montrer que $U$ ou $V$ est égal à $\{0\}$.
    \vspace{0.1cm}
    \begin{flushright}
            (\textbf{Mines-Ponts, PSI (2026) - RMS 860})
        \end{flushright}
\end{minipage}%
}
\vspace{0.4cm}

\noindent $\triangleright$ \textbf{Solution.} \\

\noindent \textbf{a)} $\bullet$ Puisque $U+V = \mathcal{L}(E)$, il existe $(p,q) \in U \times V$ tel que $\mathrm{id} = p+q$. On sait aussi que $p \circ q + q \circ p = 0$. On a $p = p \circ \mathrm{id} = p \circ (p+q) = p^2 + p \circ q$ et $p \circ q + q \circ p = p \circ q + (\mathrm{id}-p) \circ p = p \circ q + p - p^2 = 0$, d'où :
$$ \begin{cases} p = p^2 - p \circ q \\ p = p^2 + p \circ q \end{cases} $$
\noindent soit donc (par symétrie des rôles de $p$ et $q$) $p \circ q = q \circ p = 0$, et d'où :
$$ \boxed{p^2 = p \text{ et } q^2 = q} $$
\noindent d'où :
\begin{center}
\fbox{$p$ et $q$ sont des projecteurs}
\end{center}
\vspace{0.5cm}

\noindent \textbf{b)} $\bullet$ Il est clair que $\varphi$ définit une application linéaire de $V$ dans $\mathcal{L}(\ker p, E)$. Il suffit pour montrer l'injectivité de $\varphi$ de vérifier que $\ker(\varphi) = \{0\}$.

\noindent $\bullet$ Soit $v \in \ker(\varphi)$. On a donc $v_{|\ker(p)} = 0$ soit $v(x) = 0 \quad \forall x \in \ker(p)$. Il suffit donc de montrer que $v(x) = 0 \quad \forall x \in \ker(p - \mathrm{id}_E)$, car $p$ étant un projecteur, on a $E = \ker(p) \oplus \ker(p - \mathrm{id}_E)$.

\noindent $\bullet$ On sait que $p \in U$ et $v \in V$ donc $p \circ v + v \circ p = 0$. Soit $x \in \ker(p - \mathrm{id}_E)$. Alors $v(x) = v(p(x)) = -p(v(x)) \Rightarrow p(v(x)) + v(x) = 0$. En appliquant $p$ à l'égalité on obtient $p^2(v(x)) + p(v(x)) = 2p(v(x)) = 0$ d'où $p(v(x)) = 0$ et par conséquent $v(x) = 0$. On déduit que $\ker(\varphi) = \{0\}$ d'où :
\begin{center}
\fbox{$\varphi$ est injective}
\end{center}

\noindent $\bullet$ On a donc par le théorème du rang : $\dim(V) = \dim(\operatorname{Im}(\varphi))$. On peut donc tenter l'inégalité qui en résulte : $\dim(V) \leqslant \dim(\mathcal{L}(\ker p, E))$, mais par le théorème du rang appliqué à $p$, on a $\mathcal{L}(\ker p, E)$ de dimension $\dim(E) \cdot \dim(\ker p) = n(n-r) \geqslant (n-r)^2$. La borne n'est pas suffisamment restreinte.

\noindent $\bullet$ On remarque que $(n-r)^2 = \dim(\mathcal{L}(\ker p))$. Essayons de montrer que $\operatorname{Im}(\varphi) \subset \mathcal{L}(\ker p)$. Soit $v \in V$. On a donc $p \circ v + v \circ p = 0$ et donc $p \circ v_{|\ker p} + \underbrace{v \circ p_{|\ker p}}_{= 0} = 0$, d'où $\varphi(v)(x) = v_{|\ker p}(x) \in \ker(p) \quad \forall x \in \ker(p)$, et donc $\operatorname{Im}(\varphi) \subset \mathcal{L}(\ker p)$. On déduit donc que :
\begin{center}
\fbox{$\dim(V) \leqslant (n-r)^2$ avec $r = \operatorname{rg}(p)$}
\end{center}
\vspace{0.5cm}

\noindent \textbf{c)} $\bullet$ Notons $s = \operatorname{rg}(q)$. Par symétrie des rôles de $p$ et $q$, on a par le raisonnement de a) et b) : $\dim(U) \leqslant (n-s)^2$. Or, la relation $p+q = \mathrm{id}$ implique (puisque $E = \ker(p) \oplus \ker(p-\mathrm{id}) = \ker(p) \oplus \ker(q)$) que $s+r = n$, d'où $\dim(U) \leqslant (n-(n-r))^2 = r^2$ donc :
$$ \underline{\dim(U) \leqslant r^2} $$

\noindent $\bullet$ On a $\dim(\mathcal{L}(E)) = \dim(U+V) \leqslant \dim(U) + \dim(V)$ donc :
$$ \underline{n^2 \leqslant \dim(U) + \dim(V)} $$

\noindent On est en mésure de montrer que $U$ ou $V$ est égal à $\{0\}$. On a $n^2 \leqslant \dim(U) + \dim(V) \leqslant r^2 + (n-r)^2$, or on a aussi donc $n^2 - r^2 \leqslant (n-r)^2 \Leftrightarrow (n-r)(2r) \leqslant 0$. On déduit de cette inégalité que soit $r=0$, soit $n-r=s=0$, donc :
\begin{center}
\fbox{$U$ ou $V$ est égal à $\{0\}$}
\end{center}

% --- TITRE DE L'EXERCICE ---

\begin{center}

\large \textbf{4. Méthode probabiliste pour une inégalité de géometrie combinatoire}

\end{center}

\vspace{0.2cm}
% --- ENCADRÉ DE L'ÉNONCÉ ---

\noindent\framebox[\textwidth]{%

\begin{minipage}{0.94\textwidth}

\vspace{0.2cm}

Soient $E$ un espace euclidien et $(u_1, \dots, u_n)$ une famille de vecteurs unitaires de $E$. Soient $X_1, \dots, X_n$ des variables aléatoires indépendantes et de même loi uniforme sur $\{-1, 1\}$. Soit $X = \|X_1u_1 + \dots + X_nu_n\|^2$.

\vspace{0.2cm}
    \noindent \textbf{a)} Calculer $\mathbf{E}(X)$ ;
    
    \noindent \textbf{b)} Montrer qu'il existe $(\varepsilon_1, \dots, \varepsilon_n) \in \{-1, 1\}^n$ tel que $\|\varepsilon_1u_1 + \dots + \varepsilon_nu_n\| \leqslant \sqrt{n}$ .
    \vspace{0.1cm}
    \begin{flushright}
        \textbf{(Mines-Ponts, PSI (2026) - RMS 957)}
    \end{flushright}
\end{minipage}%
}

\vspace{0.4cm}
\noindent $\triangleright$ \textbf{Solution.} \\
\noindent \textbf{a)} Soit $(\Omega, \mathcal{A}, \mathbb{P})$ l'espace de probabilité en question. \\
Avant de calculer $\mathbb{E}[X]$, il faut s'assurer de son existence. On peut remarquer que $X$ ne prend qu'un nombre fini des valeurs (au nombre de $2^n$) et admet donc une espérance. Alternativement, on peut montrer que $X$ est bornée par une variable aléatoire qui admet une espérance. 
On a pour tout $\omega \in \Omega$, par l'inégalité triangulaire :
$$ |X(\omega)| = \|X_1(\omega)u_1 + \dots + X_n(\omega)u_n\|^2 \leqslant \left(|X_1| + \dots + |X_n|\right)^2(\omega) $$
\noindent Or, la variable $(|X_1| + \dots + |X_n|)^2$ est bornée, donc elle admet un moment d'ordre 1, d'où $X$ admet une espérance.

\vspace{0.3cm}
\noindent $\bullet$ On a donc :
\begin{align*}
\mathbb{E}[X] &= \mathbb{E}\left[ \left\| \sum_{i=1}^n X_i u_i \right\|^2 \right] = \mathbb{E}\left[ \sum_{i=1}^n \sum_{j=1}^n X_i X_j \langle u_i, u_j \rangle \right] \\
&= \sum_{i \neq j} \langle u_i, u_j \rangle \underbrace{\mathbb{E}[X_i X_j]}_{\substack{= \mathbb{E}[X_i]\mathbb{E}[X_j] = 0 \\ \text{par indépendance}}} + \sum_{i=1}^n \|u_i\|^2 \mathbb{E}[X_i^2]
\end{align*}
\noindent Or, puisque $X_i \sim \mathcal{U}(\{-1, 1\})$, on a $X_i^2 = 1$ $\mathbb{P}$-p.s., d'où $\mathbb{E}[X_i^2] = 1$ pour tout $i \in \llbracket 1, n \rrbracket$.

\vspace{0.3cm}
\noindent $\bullet$ On en déduit que :
\begin{align*}
\mathbb{E}[X] &= \sum_{i=1}^n \|u_i\|^2 = n \\
&\implies \boxed{\mathbb{E}[X] = n}
\end{align*}
\vspace{0.5cm}

\noindent \textbf{b)} On raisonne par l'absurde en supposant que pour tout $(\varepsilon_1, \dots, \varepsilon_n) \in \{-1, 1\}^n$, on a $\|\varepsilon_1 u_1 + \dots + \varepsilon_n u_n\| > \sqrt{n}$.

\noindent On a donc $X(\omega) = \|X_1(\omega)u_1 + \dots + X_n(\omega)u_n\|^2 > n$ pour tout $\omega \in \Omega$, d'où par croissance de l'espérance $\mathbb{E}[X] > n$, ce qui contredit directement le résultat de la question a).

\begin{center}
\fbox{\begin{tabular}{c}
Il existe $(\varepsilon_1, \dots, \varepsilon_n) \in \{-1, 1\}^n$ tel que \\
$\|\varepsilon_1 u_1 + \dots + \varepsilon_n u_n\| \leqslant \sqrt{n}$
\end{tabular}}
\end{center}
\begin{commentaire}
\textbf{Commentaire:} Cet exercice illustre l'utilisation de la \textbf{méthode probabiliste} ; une méthode robuste qui peut servir à donner une preuve non-constructive d'existence des objets mathématiques satisfaisant une certaine propriété. On s'appuie dans ce cas particulier sur le \textbf{principe des tiroirs probabiliste} qui affirme que :
\begin{center}
\fbox{Pour une variable aléatoire d'espérance $\mathbb{E}[X]$, il existe au moins une valeur de $X$ inférieure à cette espérance.}
\end{center}
\end{commentaire}

% --- TITRE DE L'EXERCICE ---
\begin{center}
    \large \textbf{5. Une équation fonctionnelle}
\end{center}
\vspace{0.2cm}

% --- ENCADRÉ DE L'ÉNONCÉ ---
\noindent\framebox[\textwidth]{%
    \begin{minipage}{0.94\textwidth}
        \vspace{0.2cm}
        Résoudre dans $\mathcal{C}^2(\mathbb{R}, \mathbb{R})$ l'équation fonctionnelle : $f(2x) = 4f(x) + 3x + 1$.
        \vspace{0.2cm}
        
        \begin{flushright}
            (\textbf{Mines-Ponts, PSI (2026) - RMS 910})
        \end{flushright}
    \end{minipage}%
}
\vspace{0.4cm}

\noindent $\triangleright$ \textbf{Solution.} \\
\noindent $\bullet$ On raisonne par analyse-synthèse.

\noindent \textbf{Analyse :} Soit $f \in \mathcal{C}^2(\mathbb{R}, \mathbb{R})$ une solution de l'équation : $f(2x) = 4f(x) + 3x + 1 \quad \forall x \in \mathbb{R}$. En dérivant deux fois l'identité, on déduit que :
$$ f''(2x) = f''(x) \quad \forall x \in \mathbb{R} $$   

\noindent $\bullet$ Une récurrence immédiate fournit donc :
$$ f''(x) = f''\left(\frac{x}{2^n}\right) \quad \forall n \in \mathbb{N} $$

\noindent Puisque $f''$ est continue, on déduit que :
$$ f''(x) = \lim_{n \to +\infty} f''\left(\frac{x}{2^n}\right) = f''(0) $$

\noindent d'où par intégration on a $f(x) = f(0) + f'(0)x + f''(0)\frac{x^2}{2}$.

\noindent Or, en exploitant l'équation et en la dérivant une fois, puis en évaluant en $0$, on obtient :
$$ \begin{array}{lcl} f(0) = 4f(0) + 3 \cdot 0 + 1 & \implies & f(0) = -\frac{1}{3} \\ f'(0) = 2f'(0) + \frac{3}{2} & \implies & f'(0) = -\frac{3}{2} \end{array} $$

\noindent L'ensemble des solutions $S$ vérifie donc :
$$ S \subset \left\{ f \in \mathcal{C}^2(\mathbb{R}, \mathbb{R}) \;\middle|\; f : x \mapsto ax^2 - \frac{3}{2}x - \frac{1}{3}, \; a \in \mathbb{R} \right\} $$

\noindent \textbf{Synthèse :} Soit $a \in \mathbb{R}$ et $f : x \mapsto ax^2 - \frac{3}{2}x - \frac{1}{3}$. On a donc :
$$ 4f(x) + 3x + 1 = f(2x) \quad \forall x \in \mathbb{R} \iff 4ax^2 - 3x - \frac{1}{3} = 4ax^2 - 6x + 3x + 1 - \frac{4}{3} $$

\noindent ce qui est vrai. pour n'importe quel $a \in \mathbb{R}$.

\noindent On déduit l'ensemble des solutions :
$$ \boxed{S = \left\{ f \in \mathcal{C}^2(\mathbb{R}, \mathbb{R}) \;\middle|\; f : x \mapsto ax^2 - \frac{3}{2}x - \frac{1}{3}, \; a \in \mathbb{R} \right\}} $$ 


\chapter{CentraleSupélec (Maths 1)}

\begin{center}

\large \textbf{1. Trace d'un endomorphisme des matrices}

\end{center}

\vspace{0.2cm}
% --- ENCADRÉ DE L'ÉNONCÉ ---

\noindent\framebox[\textwidth]{%

\begin{minipage}{0.94\textwidth}

\vspace{0.2cm}

    \noindent \textbf{a)} Montrer que : $\forall A, B \in \mathcal{M}_n(\mathbb{C}), \text{Tr}(AB) = \text{Tr}(BA)$.

    \noindent \textbf{b)} Soit $f \in \mathcal{L}(\mathcal{M}_n(\mathbb{C}), \mathcal{M}_p(\mathbb{C}))$ telle que : $\forall A, B \in \mathcal{M}_n(\mathbb{C}), f(AB) = f(A) f(B)$. Soit $\phi : M \in \mathcal{M}_n(\mathbb{C}) \mapsto \text{Tr}(f(M))$. Montrer : $\forall A, B \in \mathcal{M}_n(\mathbb{C}), \phi(AB) = \phi(BA)$ ;
    
    \noindent \textbf{c)} En déduire qu'il existe $\alpha \in \mathbb{C}$ tel que : $\forall M \in \mathcal{M}_n(\mathbb{C}), \text{Tr}(f(M)) = \alpha \text{Tr}(M)$ .
    \vspace{0.1cm}
    \begin{flushright}
        (\textbf{Centrale PSI (Maths 1) (2026) - RMS 1250})
    \end{flushright}
\end{minipage}%
}

\vspace{0.4cm}

\noindent $\triangleright$ \textbf{Solution.} \\

\noindent \textbf{a)} $\bullet$ On rappelle la définition de la trace. Soit $A \in \mathcal{M}_n(\mathbb{C})$. On définit la trace par :
$$ \boxed{\mathrm{Tr}(A) = \sum_{i=1}^n [A]_{i,i}} $$

\noindent $\bullet$ Soient $(A,B) \in \mathcal{M}_n(\mathbb{C})^2$. On a donc :
\begin{align*}
\mathrm{Tr}(AB) &= \sum_{i=1}^n [AB]_{i,i} = \sum_{i=1}^n \sum_{k=1}^n [A]_{i,k} [B]_{k,i} \\
&= \sum_{k=1}^n \sum_{i=1}^n [A]_{i,k} [B]_{k,i} = \mathrm{Tr}(BA)
\end{align*}
\noindent d'où :
$$ \boxed{\mathrm{Tr}(AB) = \mathrm{Tr}(BA)} $$
\vspace{0.5cm}

\noindent \textbf{b)} $\bullet$ On cherche à généraliser l'identité prouvée à la question a) pour $f = \mathrm{id}_{\mathcal{M}_n(\mathbb{C})}$. Soient $(A,B) \in \mathcal{M}_n(\mathbb{C})^2$. On a $\phi(AB) = \mathrm{Tr}(f(AB)) = \mathrm{Tr}(f(A)f(B))$.

\noindent Or d'après la question a), $\mathrm{Tr}(f(A)f(B)) = \mathrm{Tr}(f(B)f(A))$ puisque $(f(A), f(B)) \in \mathcal{M}_p(\mathbb{C})^2$.

\noindent On a donc $\mathrm{Tr}(f(AB)) = \mathrm{Tr}(f(B)f(A)) = \mathrm{Tr}(f(BA)) = \phi(BA)$, d'où :
$$ \boxed{\phi(AB) = \phi(BA)} $$
\vspace{0.5cm}

\noindent \textbf{c)} $\bullet$ Soit $M \in \mathcal{M}_n(\mathbb{C})$. Pour pouvoir montrer l'existence d'un $\alpha \in \mathbb{C}$ tel que $\phi(M) = \alpha \mathrm{Tr}(M)$, on a deux informations à exploiter :

\noindent $\bullet$ $\phi(AB) = \phi(BA)$ (résultat montré en b)).

\noindent $\bullet$ $\phi$ est linéaire. Il faut remarquer cela, vu que la trace est une application linéaire et que $f$ est linéaire aussi, sans quoi l'hypothèse « $f \in \mathcal{L}(\mathcal{M}_n(\mathbb{C}), \mathcal{M}_p(\mathbb{C}))$ » serait anodine.

\noindent $\bullet$ La linéarité de $\phi$ nous permet de réduire le problème sur les éléments de la base $(E_{i,j})_{1 \leqslant i,j \leqslant n}$ de $\mathcal{M}_n(\mathbb{C})$ : il suffit de vérifier que $\phi(M) = \alpha \mathrm{Tr}(M)$ pour $M = E_{i,j}$, $(i,j) \in \llbracket 1,n \rrbracket^2$.

\noindent Soit donc $(i,j) \in \llbracket 1,n \rrbracket^2$. On sait que $\mathrm{Tr}(E_{i,j}) = \delta_{i,j} = \left\{\begin{array}{ll} 1 & \text{si } i=j \\ 0 & \text{sinon} \end{array}\right.$

\noindent On a donc besoin de traiter deux cas :

\noindent $\bullet$ Si $i \neq j$, alors on veut montrer que $\phi(E_{i,j}) = 0$. On exploite la propriété $\phi(AB) = \phi(BA)$ et on rappelle la règle de calcul des produits des $(E_{i,j})$ :
$$ \boxed{E_{i,j} E_{k,l} = \delta_{j,k} E_{i,l} \quad \forall (i,j,k,l) \in \llbracket 1,n \rrbracket^4} $$

\noindent On a donc pour $A = E_{i,i}$ et $B = E_{i,j}$ :
$$ \phi(E_{i,i} E_{i,j}) = \phi(E_{i,j} E_{i,i}) $$
\noindent d'où $\phi(E_{i,j}) = 0$ puisque $E_{i,i} E_{i,j} = E_{i,j}$ et $E_{i,j} E_{i,i} = 0$.

\noindent $\bullet$ Si $i = j$, on veut montrer que $\phi(E_{i,i}) = \alpha$. Il suffit donc de montrer que $\phi(E_{i,i})$ est une constante par rapport à $i$. Soit $k \in \llbracket 1,n \rrbracket$ distinct de $i$.

\noindent On a pour $A = E_{i,k}$ et $B = E_{k,i}$ : $\phi(E_{i,k} E_{k,i}) = \phi(E_{k,i} E_{i,k})$, d'où $\phi(E_{i,i}) = \phi(E_{k,k})$. Ainsi, en posant $\alpha = \phi(E_{1,1})$, on obtient $\phi(E_{i,i}) = \alpha$.

\noindent $\bullet$ On peut conclure par linéarité de $\phi$. Pour que ce soit clair, soit $M \in \mathcal{M}_n(\mathbb{C})$. On décompose $M = \sum_{1 \leqslant i,j \leqslant n} M_{i,j} E_{i,j}$. On a donc :
\begin{align*}
\phi(M) &= \phi\left( \sum_{1 \leqslant i,j \leqslant n} M_{i,j} E_{i,j} \right) = \sum_{1 \leqslant i,j \leqslant n} M_{i,j} \phi(E_{i,j}) \\
&= \sum_{1 \leqslant i \leqslant n} M_{i,i} \alpha = \alpha \mathrm{Tr}(M)
\end{align*}
\noindent d'où :
$$ \boxed{\phi(M) = \alpha \mathrm{Tr}(M)} $$
\end{document}